\documentclass[../../project_master.tex]{subfiles}
\begin{document}
\section{From Material Damage to Device Degradation}
A useful strength of silicon radiation-damage projects is that material-level changes translate into clear device-level consequences. In this project, the preferred example is a silicon diode or detector because the physics remains transparent.

\subsection{Leakage Current}
Radiation-induced defects enhance generation-recombination activity in depleted regions, which increases reverse leakage current. A widely used reduced relation is
\begin{equation}
\Delta I = \alpha \PhiF V,
\end{equation}
where $\alpha$ is a damage coefficient and $V$ is the active volume.

\subsection{Charge Collection and Signal Loss}
If carriers recombine more rapidly after damage, then fewer of them survive long enough to be collected. This degrades charge collection efficiency, signal amplitude, and potentially detector resolution.

\subsection{Why a Diode or Detector is a Good Demonstration Case}
A diode or detector provides a direct path from defects to observables: increased leakage current, reduced diffusion length, and reduced charge collection. This makes it an ideal interview-scale example for connecting radiation transport to engineering relevance.
\end{document}
